% from ega4-8-fr.tex
\oldpage[IV-3]{50}
abéliennes [32]). Rappelons rapidement ici la notion de \emph{pro-objet} d'une
catégorie, renvoyant au chap.~V pour de plus amples développements (nous
n'utiliserons d'ailleurs pas avant le chap.~V l'interprétation qui va suivre,
et le lecteur peut donc omettre jusque-là la lecture de la fin de ce numéro).
Étant donnée une catégorie $\mathcal C$, la catégorie
$\mathbf{Pro}(\mathcal C)$ des \emph{pro-objets} de $\mathcal C$ a pour objets
les \emph{systèmes projectifs} (dans l'univers où l'on se place)
$\mathbf X=(X_\mu)_{\mu\in M}$ d'objets de $\mathcal C$ dont les ensembles
d'indices (dépendant du système projectif considéré) sont supposés préordonnés
filtrants. Étant donnés deux tels pro-objets
$\mathbf X=(X_\mu)_{\mu\in M}$,
$\mathbf X'=(X'_{\mu'})_{\mu'\in M'}$, les \emph{morphismes} de $\mathbf X$
dans $\mathbf X'$ sont par définition les éléments de l'ensemble
\[
\varprojlim_{\mu'}\left(\varinjlim_\mu
\operatorname{Hom}(X_\mu,X'_{\mu'})\right)~;
\]
la vérification du fait que l'on peut prendre ces ensembles pour ensembles de
morphismes est immédiate, la composition de systèmes de morphismes
$u_{\mu'}^\mu:X_\mu\to X'_{\mu'}$,
$u_{\mu''}^{\prime\mu'}:X'_{\mu'}\to X''_{\mu''}$ qui sont inductifs en
l'indice supérieur et projectifs en l'indice inférieur, se faisant
«~argument par argument~», autrement dit en considérant le système des
$u_{\mu''}^{\prime\mu'}\circ u_{\mu'}^\mu$.
\end{env}

\begin{env}[8.13.4]
\label{IV.8.13.4-fr}
Considérons alors un préschéma quasi-compact et quasi-séparé $S$, et désignons
par $\mathcal C$ la sous-catégorie pleine de la catégorie
$(\mathbf{Sch})_{/S}$ des $S$-préschémas, formée des $S$-préschémas $X$ ayant
la propriété suivante~: le morphisme structural $X\to S$ se factorise en
\[
X\xrightarrow{g}X_0\xrightarrow{f}S,
\]
où $g:X\to X_0$ est affine et $f:X_0\to S$ de présentation finie~; nous
dirons pour abréger que les préschémas de $\mathcal C$ sont
\emph{essentiellement affines au-dessus de $S$}.

Considérons d'autre part la sous-catégorie pleine $\mathcal C'_0$ de
$(\mathbf{Sch})_{/S}$ formée des $S$-préschémas \emph{de présentation finie},
et la catégorie $\mathbf{Pro}(\mathcal C'_0)$ des pro-objets de
$\mathcal C'_0$. Nous dirons qu'un objet
$\mathbf X=(X_\mu)_{\mu\in M}$ de $\mathbf{Pro}(\mathcal C'_0)$ est
\emph{essentiellement affine} s'il existe un $\gamma\in M$ tel que pour tout
$\mu\geqslant\gamma$, le morphisme de transition
$v_{\gamma\mu}:X_\mu\to X_\gamma$ soit affine (ce qui entraîne que pour
$\gamma\leqslant\mu\leqslant\nu$, $v_{\mu\nu}:X_\nu\to X_\mu$ est affine).
On notera qu'un objet de $\mathbf{Pro}(\mathcal C'_0)$ \emph{isomorphe} à un
objet essentiellement affine \emph{n'est pas} nécessairement essentiellement
affine lui-même.

Nous désignerons par $\mathcal C'$ la sous-catégorie pleine de
$\mathbf{Pro}(\mathcal C'_0)$ formée des pro-objets de $\mathcal C'_0$
essentiellement affines.

Cela étant, il résulte de (8.2.2) et (8.2.3) que pour tout objet
$\mathbf X=(X_\mu)_{\mu\in M}$ de $\mathcal C'$, le $S$-préschéma
$X=\varprojlim_\mu X_\mu$ existe~; en outre, comme, pour $\mu$ assez grand,
le morphisme $X\to X_\mu$ est affine (8.2.2), $X$ est
\emph{essentiellement affine au-dessus de $S$} par définition. Posons
$X=L(\mathbf X)$~; montrons qu'on a ainsi défini un \emph{foncteur canonique}
\begin{equation}
\label{IV.8.13.4.1-fr}
L:\mathcal C'\longrightarrow\mathcal C. \tag{8.13.4.1}
\end{equation}

On a en effet, pour deux objets $\mathbf X=(X_\mu)$,
$\mathbf X'=(X'_{\mu'})$ de $\mathcal C'$, une application canonique pour
chaque $\mu'$
\[
\varinjlim_\mu\operatorname{Hom}_S(X_\mu,X'_{\mu'})
\longrightarrow
\operatorname{Hom}_S(\varprojlim_\mu X_\mu,X'_{\mu'})
\]
définie dans (8.13.1.1), et d'autre part, par définition de la limite
projective, une bijection canonique
\[
\varprojlim_{\mu'}
\operatorname{Hom}_S(\varprojlim_\mu X_\mu,X'_{\mu'})
\xrightarrow{\sim}
\operatorname{Hom}_S(\varprojlim_\mu X_\mu,
\varprojlim_{\mu'}X'_{\mu'}),
\]
